\label{chap:background}

\section{Posizioni, mosse e problemi di matto}

Una posizione scacchistica è descritta nel progetto mediante la notazione FEN
(\textit{Forsyth--Edwards Notation}), che codifica disposizione dei pezzi, lato al
tratto, diritti di arrocco, eventuale casa en passant e contatori di mossa. Le mosse
sono espresse in formato UCI: la stringa identifica casa di origine, casa di
destinazione ed eventuale promozione, per esempio \texttt{e2e4} o
\texttt{e7e8q}. Il formato è compatto e può essere verificato automaticamente
rispetto alle regole della posizione.

Un problema \textit{mate-in-N} richiede una prima mossa che garantisca il matto
entro il numero dichiarato di mosse del solutore, contro ogni difesa ammessa. Nel
dataset Lichess la FEN precede una mossa di preparazione dell'avversario: la
posizione mostrata al solutore si ottiene applicando la prima mossa della sequenza,
mentre la seconda è la prima mossa della soluzione
\cite{lichessOpenDatabase}. GraphMate apprende questa singola mossa target; non
genera autonomamente l'intera variante fino al matto.

\section{Reti neurali su grafi e attenzione}

Una rete neurale a grafo aggiorna la rappresentazione di ciascun nodo aggregando
informazioni provenienti dai nodi collegati. Ripetendo l'operazione su più strati,
ogni nodo incorpora un contesto strutturale progressivamente più ampio. Il
\textit{pooling} riassume poi le rappresentazioni locali in un vettore del grafo,
utilizzabile per una decisione relativa all'intera posizione. Questa impostazione
generalizza l'idea delle convoluzioni su grafi introdotta, tra gli altri, da
Kipf e Welling \cite{kipf2017semi}.

Le \textit{Graph Attention Networks} (GAT) assegnano pesi appresi ai vicini invece
di combinarli tutti nello stesso modo \cite{velickovic2018graph}. Intuitivamente,
quando il modello aggiorna la rappresentazione di una casa, può attribuire maggiore
rilevanza alle relazioni più informative per quella posizione. L'attenzione è
calcolata separatamente da più teste e i risultati vengono combinati; ciò consente
di apprendere prospettive differenti sulle stesse relazioni. In GraphMate gli archi
portano inoltre attributi scacchistici, cosicche il messaggio dipende sia dai nodi
coinvolti sia dal tipo di relazione. Il Capitolo~4 specifica le feature effettive e
il Capitolo~5 descrive l'encoder adottato.

\section{Informazione temporale e rapporto con TimeGNN}

Le reti temporali su grafi introducono variabili di tempo nella costruzione o
nell'elaborazione delle relazioni. TimeGNN, per esempio, studia rappresentazioni
temporali dinamiche per serie multivariate \cite{xu2023timegnn}. La libreria
TimeGNN presente nel progetto ha fornito riferimenti architetturali e tecniche di
implementazione basate su PyTorch Geometric, come layer di attenzione, gestione
degli attributi degli archi e organizzazione dell'addestramento.

GraphMate non esegue però il framework TimeGNN completo e non tratta la scacchiera
come un event log dinamico. I suoi grafi rappresentano singole posizioni; nel
Modello B il timing sintetico è codificato da un piccolo encoder e concatenato al
contesto globale usato per valutare le mosse candidate. Si tratta quindi di un
adattamento di idee e tecniche, progettato per un compito graph-level differente.

\section{Modelli linguistici come baseline}

Un modello linguistico può ricevere la FEN in un prompt e generare testo contenente
una mossa. Diversamente dai modelli GraphMate, non riceve in ingresso il grafo
scacchistico costruito dal progetto e non applica internamente il generatore di
mosse legali della pipeline GNN. Per ottenere un confronto verificabile, il
protocollo congela prompt, parametri di generazione e parser e richiede una singola
mossa UCI. La valutazione separa output non interpretabili, mosse illegali, mosse
legali errate e risposte corrette; i dettagli sono riportati nel Capitolo~6 e
nell'Appendice~\ref{app:llm-protocol}.
